\section{Reproducibility and Artifact Map}

\subsection{Commands and environment}
\begin{draftpoints}
  \item project environment is defined by \artifact{pyproject.toml}, \artifact{uv.lock}, and the source tree under \artifact{src/}; paper-shaped configurations are stored in \artifact{config/paper/}
  \item Main paper configurations can be executed with the following command:
  \begin{quote}\footnotesize\ttfamily
  python scripts/run\_paper\_config.py\\
  \hspace*{1em}--config config/paper/figure4.yaml
  \end{quote}
  \item compact suite runner is \texttt{python scripts/run\_paper\_experiments.py}
  \item Portable data can be regenerated with the following command; report plots have dedicated scripts under \artifact{scripts/}:
  \begin{quote}\footnotesize\ttfamily
  python scripts/export\_replication\_figure\_data.py
  \end{quote}
  \item full protocol is compiled from \artifact{docs/protocol/main.tex} with two LaTeX passes around BibTeX
\end{draftpoints}

\subsection{Portable result tables and figures}
\begin{table}[htbp]
\centering
\caption{Protocol figures and their portable source tables. Paths are relative to \artifact{docs/figures/replication/data/}.}
\label{tab:figure-data-map}
\scriptsize
\begin{tabularx}{\textwidth}{p{4.3cm}Y}
\toprule
Figure & Portable data \\
\midrule
MNIST performance & \artifact{mnist_condition_metrics.csv}; \artifact{mnist_seed_metrics.csv} \\
Claim-direction ratios & \artifact{mnist_claim_ratios.csv} \\
Per-class MNIST & \artifact{mnist_per_class_summary.csv}; \artifact{mnist_per_class_seed.csv} \\
Growth mechanism & \artifact{mnist_architecture.csv}; MNIST condition and seed tables \\
SD-19 feasibility & \artifact{sd19_plotted_summary.csv}; \artifact{sd19_seed_metrics.csv}; \artifact{sd19_per_class_seed.csv} \\
Cross-dataset ratios & \artifact{cross_dataset_ratios.csv} \\
Training resources & \artifact{training_resource_conditions.csv}; \artifact{training_resource_ratios.csv} \\
Predictive coding & \artifact{predictive_coding_aggregate.csv}; \artifact{predictive_coding_seed_metrics.csv} \\
Post-replication optimization & three \artifact{post_replication_optimization_*.csv} tables \\
Complete bundle & \artifact{replication_figure_base_data.json} \\
\bottomrule
\end{tabularx}
\end{table}

\subsection{Preserved and missing artifacts}
\begin{draftpoints}
  \item Preserved evidence includes the full implementation, experiment configurations, tests, reports, nine generated replication figures, tidy CSV tables, and an 8.46 MB consolidated JSON bundle
  \item CSV exports contain 60 MNIST seed rows, 600 MNIST per-class seed rows, six SD-19 seed rows, 96 SD-19 per-class rows, predictive-coding results, optimization screens, and resource measurements
  \item Only one local MLflow diagnostic run remains in the workspace; complete raw logs and checkpoints for the reported confirmation are absent
  \item Raw MNIST and SD-19 datasets are not versioned and must be obtained through the configured data modules before rerunning experiments
  \item Missing checkpoints prevent retrospective generation of the required matched qualitative grids; reruns should treat checkpoint retention and fixed-sample manifests as mandatory outputs
\end{draftpoints}

\subsection{Recommended rerun outputs}
\begin{draftpoints}
  \item Save the resolved configuration, code revision, environment lock, hardware description, and immutable dataset split for every seed
  \item Save base, per-class, and final checkpoints for every confirmatory condition
  \item Export the uniform trajectory schema in Table~\ref{tab:trajectory-schema} and qualitative manifest in Table~\ref{tab:qualitative-manifest}
  \item Generate fixed-sample final and curriculum-step reconstruction grids immediately after each run so visual evidence does not depend on later checkpoint availability
\end{draftpoints}
